\clearpage\thispagestyle{empty}
\pdfbookmark[0]{Quellen zum Juno-Bericht vom Mai 1805}{juno180505-sources}
\begin{center}\textbf{Quellen zum Juno-Bericht vom Mai 1805}\end{center}
\noindent Diese KI-Übertragung in heutiges Deutsch folgt dem Auszug aus Franz Xaver von Zachs Bericht in \href{https://gdz.sub.uni-goettingen.de/id/PPN236007467}{Gauß’ \textit{Werken}, Band~VI} (Göttingen, 1874), S.~263--265. Der Bericht erschien als Beitrag XLII, \textit{Fortgesetzte Nachrichten über den neuen Haupt-Planeten Juno}, in der \textit{Monatlichen Correspondenz}, Band~XI, Mai 1805. Die Beobachtungen von Giuseppe Piazzi in Palermo stehen auf \href{https://www.digitale-sammlungen.de/en/view/bsb10538603?page=489}{S.~475}, der Vergleich und die fünften Elemente von Carl Friedrich Gauß auf \href{https://www.digitale-sammlungen.de/en/view/bsb10538603?page=490}{S.~476}. Auf \href{https://www.digitale-sammlungen.de/en/view/bsb10538603?page=491}{S.~477} stehen Gauß’ Beobachtung vom 20. Februar und der Anfang der Ephemeride. Die Zeitschrift berichtet auch über Beobachtungen von Barnaba Oriani und Wilhelm Olbers. Sie nennt Johann Karl Burckhardt als Urheber der Bahnelemente, die sie auf \href{https://www.digitale-sammlungen.de/en/view/bsb10538603?page=496}{S.~482} nach dem \textit{Moniteur} vom 28. Dezember 1804 anführt.
\par\smallskip\noindent Piazzi beschreibt den Vollkreis in Buch~II seiner \href{https://www.e-rara.ch/zut/content/titleinfo/11884975}{\textit{Della specola astronomica de’ regi studi di Palermo}} (Palermo, 1792). Die in der Werkausgabe angeführten Elemente zur Berliner Zeit stehen auf \href{https://www.digitale-sammlungen.de/en/view/bsb10538337?page=280}{S.~270} in Johann Elert Bodes \href{https://www.digitale-sammlungen.de/en/view/bsb10538337?page=7}{\textit{Astronomischem Jahrbuch für das Jahr 1808}} (Berlin, 1805).
\par\smallskip\noindent Übertragung und Satz wurden mit KI erstellt. Die Digitalisate von Zeitschrift und Jahrbuch stammen von der Bayerischen Staatsbibliothek.
\clearpage
